%%% Please fill in basic information on your thesis, which will be automatically
%%% inserted at the right places. You need to replace \xxx{...} by real data.

% Type of your thesis:
%	"bc" for Bachelor's
%	"mgr" for Master's
%	"phd" for PhD
%	"rig" for rigorosum
\def\ThesisType{mgr}

% Language of your study programme:
%	"cs" for Czech
%	"en" for English
\def\StudyLanguage{cs}

% Thesis title in English (exactly as in the official assignment)
% (Note: \xxx is a "ToDo label" which makes the unfilled visible. Remove it.)
\def\ThesisTitle{Defining robot trajectory using augmented reality}

% Author of the thesis (you)
\def\ThesisAuthor{Jakub Medek}

% Year when the thesis is submitted
\def\YearSubmitted{2026}

% Name of the department or institute, where the work was officially assigned
% (according to the Organizational Structure of MFF UK in English,
% see https://www.mff.cuni.cz/en/faculty/organizational-structure,
% or a full name of a department outside MFF)
\def\Department{Department of Theoretical Computer Science and Mathematical Logic}

% Is it a department (katedra), or an institute (ústav)?
\def\DeptType{Department}

% Thesis supervisor: name, surname and titles
\def\Supervisor{RNDr. David Obdržálek, Ph.D.}

% Supervisor's department (again according to Organizational structure of MFF)
\def\SupervisorsDepartment{Department of Theoretical Computer Science and Mathematical Logic}

% Study programme (does not apply to rigorosum theses)
\def\StudyProgramme{Computer Science - Artificial Intelligence}

% An optional dedication: you can thank whomever you wish (your supervisor,
% consultant, who provided you with tea and pizza, etc.)
\def\Dedication{%
I would like to thank my supervisor, RNDr. David Obdržálek, Ph.D., for his guidance and help during the creation of this work.
I would also like to thank my family for their continued support and encouragement throughout my studies.
}

% Abstract (recommended length around 80-200 words; this is not a copy of your thesis assignment!)
\def\Abstract{%
This work investigates whether collaborative-robot tasks that change frequently can be authored more directly by specifying spatial intent in the real workspace instead of editing low-level robot programs through a teach pendant or offline workflow.
It presents a modular system for task-oriented in-workspace spatial authoring: a shared runtime provides common services for sensing, authoring, confirmation, execution, and persistence, while task-specific use cases interpret captured poses and trajectories as robot motions and tool actions.
The system was implemented on a Kassow Robots collaborative robot using stereo-vision workspace sensing and a web-based interface.
As part of the system, this work also designed and built a custom tracked 6DoF pen that provides spatial input and button-based intent events.
Two representative use cases were demonstrated on the real robot: pick-and-place and seam welding.
The system demonstrates that this approach can support task-oriented workflows, with seam welding showing the clearest practical fit because authored trajectory input is used to identify and parameterize the weld seam, while pick-and-place shows that authored intent can be combined with refreshed scene information at execution time.
}

% 3 to 5 keywords (recommended) separated by \sep
% Keywords are useful for indexing and searching for the theses by topic.
\def\ThesisKeywords{%
spatial authoring\sep
collaborative robot\sep
task-oriented robot programming\sep
augmented reality\sep
modular software architecture%
}

% If any of your metadata strings contains TeX macros, you need to provide
% a plain-text version for use in XMP metadata embedded in the output PDF file.
% If you are not sure, check the generated thesis.xmpdata file.
\def\ThesisAuthorXMP{\ThesisAuthor}
\def\ThesisTitleXMP{\ThesisTitle}
\def\ThesisKeywordsXMP{\ThesisKeywords}
\def\AbstractXMP{\Abstract}

% If your abstracts are long and do not fit in the infopage, you can make the
% fonts a bit smaller by this setting. (Also, you should try to compress your abstract more.)
\def\InfoPageFont{}
%\def\InfoPageFont{\small}  % uncomment to decrease font size

% If you are studing in a Czech programme, you also need to provide metadata in Czech:
% (in English programmes, this is not used anywhere)

\def\ThesisTitleCS{Zadávání trajektorie robota pomocí rozšířené reality}
\def\DepartmentCS{Katedra teoretické informatiky a matematické logiky}
\def\DeptTypeCS{Katedra}
\def\SupervisorsDepartmentCS{Katedra teoretické informatiky a matematické logiky}
\def\StudyProgrammeCS{Informatika - Umělá inteligence}

\def\ThesisKeywordsCS{%
prostorové zadávání\sep
kolaborativní robot\sep
úlohově orientované programování robotů\sep
rozšířená realita\sep
modulární softwarová architektura%
}

\def\AbstractCS{%
Tato práce zkoumá, zda lze často se měnící úlohy pro~kolaborativní roboty zadávat příměji specifikací prostorového záměru v~reálném pracovním prostoru namísto úprav nízkoúrovňových programů robota pomocí teach pendantu nebo offline postupů.
Práce představuje modulární systém pro~úlohově orientované prostorové zadávání v~pracovním prostoru: sdílená běhová vrstva poskytuje obecné služby pro~snímání scény, zadávání pozic a~trajektorií, potvrzování, spouštění a~ukládání, zatímco úlohově specifické moduly interpretují zachycené pozice a~trajektorie jako pohyby robota a~akce nástroje.
Systém byl implementován na~kolaborativním robotu Kassow Robots s~využitím stereoskopického snímání pracovního prostoru a~webového rozhraní.
V~rámci vývoje systému bylo také navrženo a~zkonstruováno vlastní sledované pero se šesti stupni volnosti (6DoF), které poskytuje prostorový vstup a~slouží k~zachycení záměru uživatele prostřednictvím tlačítek.
Na~reálném robotu byly demonstrovány dvě reprezentativní úlohy: pick-and-place a~svařování.
Výsledky demonstrují, že tento přístup může podporovat úlohově orientované pracovní postupy.
Nejpraktičtější uplatnění ukázalo svařování, kde zadaná trajektorie slouží přímo k~lokalizaci a~parametrizaci svaru. Úloha pick-and-place naopak dokládá, že uživatelsky zadaný záměr lze při vykonání programu kombinovat s~aktualizovanými informacemi o~scéně.
}
